\section{Unitarizing every simple module}
\label{sec:unitarity}

Theorem~\ref{thm:positive-fusion} closes the unitary modules under
algebraic fusion. We now use this closure to reach every simple module.

\begin{theorem}\label{thm:strong-unitarity}
Let $V$ be a simple unitary strongly rational vertex operator algebra.
Every simple grading-restricted $V$-module admits a positive definite
invariant Hermitian form. Consequently, every grading-restricted
finite-length $V$-module is unitarizable, and $V$ is strongly unitary.
\end{theorem}

Here is how a module with no given positive form enters the unitary
theory. Fix a simple $V$-module $P$, and set
\[
 D=V\otimes V,\qquad
 U=D^{\langle\sigma\rangle},\qquad R=P\otimes P^\vee,
\]
where $\sigma$ exchanges the factors. The transposition-twisted vacuum
module $T=T_\sigma(V)$ is unitary over $U$. The permutation fusion
formula below produces a nonzero twisted intertwiner of type
$\binom{T}{R\ T}$. Its charge map is injective because $R$ is simple
over $D$. Consequently it detects every simple constituent of $R|_U$,
placing that constituent in the positive fusion product
$T\boxtimes_U T^\vee$. This gives a positive $U$-invariant form $h_1$
on $R$.

Restriction to $U$ controls only the flip-invariant fields. To recover
the adjoint identity for the remaining fields, write $D=U\oplus J$,
where $J=\{d\in D:\sigma(d)=-d\}$ is the sign constituent. Its action identifies
$J\boxtimes_U R$ with $R$. Transferring the canonical fusion form along
that action gives a second form $h_2$. The normalized transport identity
interchanges $h_1$ and $h_2$ when taking the adjoint of a $J$-field, so
$h_1+h_2$ is invariant under all of $D$. Its restriction to
$P\otimes\mathbb Cq$, for any nonzero $q\in P^\vee$, then unitarizes
$P$. We prove the algebraic permutation formula first, without using
unitarity of $P$.

\subsection{The algebraic permutation product}

\begin{lemma}[The required permutation product before positivity]
\label{lem:permutation-before-positivity}
Let \(V\) be simple and strongly rational, let \(D=V\otimes V\),
and let \(\sigma\) exchange the factors. Write \(T_\sigma(N)\) for
the Barron--Dong--Mason transposition-twisted module associated to an
ordinary \(V\)-module \(N\) \cite[Theorem~3.9]{BDM}.
The fusion products below are taken among grading-restricted
ordinary or \(\sigma\)-twisted \(D\)-modules. For every ordinary
\(V\)-module \(P\),
\[
 (P\otimes P^\vee)\boxtimes_D T_\sigma(V)
 \cong T_\sigma(P\boxtimes_V P^\vee).
\]
No unitarity or positivity of the lowest conformal weight of \(P\)
is assumed.
\end{lemma}

\begin{proof}
Put \(U=D^{\langle\sigma\rangle}\).  Cyclic-orbifold regularity
\cite[Theorem~5.24 and Corollary~5.25]{CarnahanMiyamoto} makes
\(U\) strongly rational.  Let \(\mathcal C\) be its category of
grading-restricted generalized modules, with its HLZ vertex tensor
category structure, and write \(c_{X,Y}\) for its braiding.
We use \(L(n)=L_n\) for Virasoro modes in the coordinate formulas below.
Ordinary \(D\)-modules and the BDM modules \(T_\sigma(N)\)
restrict to objects of \(\mathcal C\):
\[
 L^{T_\sigma(N)}(0)=\tfrac12L^N(0)+\tfrac{c_V}{16}.
\]
This requires boundedness below and finite-dimensional weight spaces,
not nonnegative module weights.

Regard \(D\) as a commutative algebra object in \(\mathcal C\), and
write \(\star\) for the relative product on
\(\operatorname{Rep}_{\mathcal C}(D)\).  This product is the cokernel
imposing equality of the two \(D\)-actions.  Right exactness allows
the ambient associator and its inverse to descend through these
quotients, giving an associator for \(\star\)
\cite[Theorems~2.51--2.53]{CKMExtensions}.
If \(\mu_Q\) denotes the action on an object \(Q\), the categorical
\(g\)-twisted condition is
\[
 \mu_Q(g\boxtimes 1_Q)c_{Q,D}c_{D,Q}=\mu_Q .
\]
For the canonical quotient \(q:X\boxtimes_U Y\to Q=X\star Y\) of
objects in the \(g\)- and \(h\)-sectors, the target lies in the \(gh\)-sector:
precompose this identity with \(1_D\boxtimes q\),
use the hexagon and balancing, and then apply the two input identities.
Right exactness makes that precomposition surjective
\cite[Proposition~2.4 and Corollary~2.5]{McRaeTwisted}.
The field attached to \(\mu_Q\) has the vacuum, truncation, derivative,
and grading properties; the displayed monodromy condition gives
the equivariance for its prescribed twist.  Analytic associativity reconstructs its twisted
Jacobi identity
\cite[Propositions~4.11 and~4.13, Remark~4.12 and Theorem~4.14]{McRaeTwisted}.
Thus these are actual twisted modules.  Only the known \(1\)- and
\(\sigma\)-sectors are needed, not an exhaustive classification of
algebra modules.

We spell out the relevant intertwiner dictionary.
The two categorical action compatibilities are precisely analytic
associativity and skew-associativity for the corresponding
\(U\)-intertwiner \(\mathcal I\)
\cite[Theorem~3.44 and Proposition~3.50]{CKMExtensions}.
Suppose its charge \(X\) is ordinary and its source and target are
\(\sigma\)-twisted.  For \(\sigma a=(-1)^j a\) with \(j\in\{0,1\}\), the matrix
coefficients of
\[
 \begin{gathered}
 Y_{\rm target}(a,z_1)\mathcal I(x,z_2)y,\\
 \mathcal I(x,z_2)Y_{\rm source}(a,z_1)y,\\
 \mathcal I(Y_X(a,z_1-z_2)x,z_2)y
 \end{gathered}
\]
are branches of one analytic function \(g(z_1,z_2)\).
For fixed nonzero \(z_2\), multiplication by \(z_1^{j/2}\)
cancels monodromy about \(0\).  The iterate has integer powers
about \(z_1=z_2\), since \(X\) is ordinary.  These two local loops
generate the punctured-plane fundamental group, so
\(z_1^{j/2}g(z_1,z_2)\) is single-valued.  Source truncation gives a pole of finite order at \(0\), and
ordinary-charge truncation gives one at \(z_1=z_2\). At infinity,
the finite weight support of the restricted-dual vector and the lower
bound on intermediate target weights give polynomial growth. Thus the
only possible poles are at \(0,z_2,\infty\).
To pass from this fixed-$z_2$ rationality to a two-variable formal
identity, take homogeneous
inputs of weights $h_x,h_y$, a state $a$ of weight $d$, and a
restricted-dual vector of weight $H$. On compatible branches, the
$L(0)$-bracket gives
\[
 z_1^{j/2}g(z_1,z_2)
 =z_2^{H-d-h_x-h_y+j/2}
   (z_1/z_2)^{j/2}g(z_1/z_2,1).
\]
The function $u\mapsto u^{j/2}g(u,1)$ is rational, with poles only at $0,1,\infty$;
the prefactor is a rational power. Its three expansions therefore give
the twisted Jacobi identity. The converse follows from formal duality
\cite[Proposition~2.17]{DLXY}.
Semisimple \(L(0)\) and the \(L(0)\)-bracket remove logarithmic terms.
Here the factor is \(z_1^{j/2}\), as in Proposition~2.17; this corrects
the interchange of the two exponents in the displayed cancellation
factor in the proof of \cite[Theorem~3.5]{DLXY}.
Consequently \(\star\) represents exactly the needed twisted
intertwining operators.  On ordinary modules it is ordinary
\(D\)-fusion \cite[Theorem~3.54 and Corollary~3.55]{CKMExtensions}.

For an integer $k\geq1$, define the coefficients $a_j$ (depending on
$k$) by the formal coordinate identity
\begin{equation}
 \exp\!\left(-\sum_{j\geq1}a_jx^{j+1}\frac{d}{dx}\right)x
       =\frac{(1+x)^k-1}{k},
 \label{eq:permutation-delta-coordinate}
\end{equation}
and set
\begin{equation}
 \Delta_k(z)=\exp\!\left(\sum_{j\geq1}a_jz^{-j/k}L(j)\right)
                 k^{-L(0)}z^{(1/k-1)L(0)}.
 \label{eq:permutation-delta-operator}
\end{equation}
These are the coordinate operators of
\cite[Section~5.1.1, equations~(5.3)--(5.5)]{DLXY}, with the chosen
branches of the indicated powers.

Write $I_V$ and $I_D$ for the indicated spaces of ordinary and twisted
intertwining operators, respectively. For ordinary $V$-modules $M,N$,
the explicit mutually inverse maps of
\cite[Theorems~5.5, 5.9 and Corollary~5.10]{DLXY} give
\[
 \begin{gathered}
 I_V\binom{W}{M\ N}
 \cong I_D\binom{T_\sigma(W)}{M\otimes V\ T_\sigma(N)},\\
 \overline{\mathcal Y}(m\otimes\mathbf1,z)
 =\mathcal Y(\Delta_2(z)m,z^{1/2}).
 \end{gathered}
\]
Together with the BDM equivalence and the ordinary fusion universal
property, this identifies the representing object:
\begin{equation}
 (M\otimes V)\star T_\sigma(N)
 \cong T_\sigma(M\boxtimes_V N).
 \label{eq:permutation-one-slot}
\end{equation}
This is the universal-property proof of
\cite[Theorem~6.3]{DLXY}, applied only to grading-restricted targets.

For completeness, the finite-expansion step in that proof does not
require integral module weights.  If \(L(0)m=hm\), write
\[
 e^{-\sum_{j\ge1}a_jz^{-j}L(j)}m=\sum_s z^{-s}m_s,
 \qquad L(0)m_s=(h-s)m_s.
\]
The sum is finite by boundedness below, and
\[
 \Delta_k(z^k)^{-1}m
 =z^{(k-1)h}\sum_s k^{h-s}z^{-ks}m_s.
\]
In the weak-associativity substitution this fractional monomial
depends only on \(z_2\) and multiplies both sides.
The corresponding expansion of a state \(u\in V\) has integral
powers.  Taking a maximum over the finitely many coefficient pairs
and absorbing their finitely many negative integral powers supplies
the common exponent used in the proof of
\cite[Theorem~5.9]{DLXY}.

To move the slot, pull back all three actions in an intertwiner by
\(\sigma\).  This preserves twisted Jacobi and its universal property.
For every ordinary \(\sigma\)-twisted object \(Q\),
\[
 E_Q=e^{-2\pi iL^Q(0)},\qquad
 E_Qa_nE_Q^{-1}=(-1)^j a_n
 \quad
 \bigl(\sigma a=(-1)^j a,\ n\in j/2+\mathbb Z\bigr).
\]
Thus \(E_Q:Q\longrightarrow Q^\sigma\) is an invertible module map.
The charge pullback of \(M\otimes V\) is \(V\otimes M\).
Applying simultaneous pullback to \eqref{eq:permutation-one-slot} gives
\[
 (V\otimes M)\star T_\sigma(N)
 \cong T_\sigma(M\boxtimes_V N).
\]
This does not require \(M\otimes V\) itself to be flip-invariant.

Finally, put \(X=P\otimes V\), \(Y=V\otimes P^\vee\).
Ordinary tensor-product factorization and the vacuum unit give
\(X\star Y\cong P\otimes P^\vee\).  Its factorized intertwiner is
\[
 (p\otimes v,\ u\otimes q)\longmapsto
 e^{zL_P(-1)}Y_P(u,-z)p\otimes Y_{P^\vee}(v,z)q,
\]
where \(u\) belongs to the first factor and \(v\) to the second.
The relative associator and the two slot formulas now yield
\[
 \begin{aligned}
 (P\otimes P^\vee)\star T_\sigma(V)
 &\cong (X\star Y)\star T_\sigma(V)\\
 &\cong X\star\bigl(Y\star T_\sigma(V)\bigr)\\
 &\cong (P\otimes V)\star T_\sigma(P^\vee)\\
 &\cong T_\sigma(P\boxtimes_V P^\vee).
 \end{aligned}
\]
The dictionary permits writing \(\boxtimes_D\) for these products.
If \(V=\mathbb C\), the flip is the identity and the statement is
immediate.
\end{proof}

\subsection{From detection to a positive invariant form}

\begin{proof}[Proof of Theorem~\ref{thm:strong-unitarity}]
For $V=\mathbb C$, the only simple module is the scalar module, so the
assertion is immediate. Assume henceforth that $V\ne\mathbb C$, and fix
an arbitrary simple $V$-module $P$. Its algebraic contragredient is
denoted by $P^\vee$; neither module is assumed unitary.

\smallskip
\noindent\emph{The fixed-point algebra and a unitary twisted module.}
Set
\[
 D=V\otimes V,\qquad
 \sigma(v\otimes w)=w\otimes v,\qquad
 U=D^{\langle\sigma\rangle},\qquad J=D^{-\sigma}.
\]
Thus $D=U\oplus J$ as a $U$-module. Tensor products preserve simplicity,
regularity, CFT type and self-contragredience. In particular, $D$ is a
simple unitary strongly rational VOA. The flip preserves its unitary
form and commutes with its PCT operator $\Theta^D$. The cyclic
fixed-point theorem of Carnahan--Miyamoto
\cite[Corollary~5.25]{CarnahanMiyamoto} applies to $D$: it is simple,
nonnegatively graded and regular, with a nonsingular invariant bilinear
form. It follows that $U$ is regular. CFT type and the unitary structure
are inherited, and the inherited form and PCT also give
self-contragredience. Hence $U$ is simple unitary strongly rational,
and Theorem~\ref{thm:positive-fusion} applies to $U$.

Let $T=T_\sigma(V)$ be the transposition-twisted $D$-module obtained from
the vacuum $V$ by the Barron--Dong--Mason construction
\cite[Theorem~3.9]{BDM}. This construction preserves the unitary form
\cite[Lemma~10.12]{GuiPositivity}, so $T|_U$ is unitary. Its conformal
energy is
\[
 L_0^T=\tfrac12 L_0^V+\tfrac{c}{16},
\]
where $c$ is the central charge of $V$. In particular, it has
finite-dimensional energy spaces and a lower bound. Put
\[
 R=P\otimes P^\vee.
\]
The tensor-product module theorem makes $R$ simple over $D$; it too has
finite-dimensional energy spaces and a lower bound.

The restrictions $T|_U$ and $R|_U$ have the same conformal vector and
energy spaces as before restriction. They are grading-restricted
$U$-modules of finite length. Indeed, regularity decomposes them into
simple ordinary $U$-modules. There are finitely many simple types, and
infinitely many copies of one type would make the energy space at its
lowest weight infinite-dimensional. All fusion products over $U$ used
below therefore belong to its finite semisimple module category.

\smallskip
\noindent\emph{A positive form on the restriction of $R$.}
Here $\boxtimes_D$ denotes the product representing twisted
intertwining operators among grading-restricted ordinary or
$\sigma$-twisted $D$-modules, as constructed in
Lemma~\ref{lem:permutation-before-positivity}. That lemma gives
\begin{equation}\label{eq:permutation-fusion-detection}
 R\boxtimes_D T_\sigma(V)
 \cong T_\sigma(P\boxtimes_V P^\vee).
\end{equation}
The ordinary labels $P$ and $P^\vee$ are still not assumed unitary.
By rigidity, $P\boxtimes_V P^\vee$ has a nonzero evaluation onto $V$.
Applying $T_\sigma$ to this morphism and using
\eqref{eq:permutation-fusion-detection} yields a nonzero twisted
intertwining operator $I$: the universal intertwiner's coefficients
generate its target, so composition with a nonzero target morphism
cannot vanish. Its type is
\[
 \binom{T}{R\ T}.
\]

We need this operator to detect every simple summand of $R|_U$.
For two weak $\sigma$-twisted $D$-modules, the space of generalized
intertwining operators between them is a weak $D$-module
\cite[Theorem~4.6]{DLXY}; write $\mathcal H(T,T)$ for this space with
source and target $T$. Moreover,
\cite[Proposition~4.8]{DLXY} identifies the charge map
\[
 f_I:R\longrightarrow\mathcal H(T,T),\qquad p\longmapsto I(p,x),
\]
as a homomorphism of weak $D$-modules. Its kernel is therefore a
$D$-submodule. Since $R$ is simple and $I$ is nonzero, $f_I$ is
injective.

Restricting the observable states in the twisted Jacobi identity to
$U$ gives the ordinary Jacobi identity, with the same derivative
identity and lower truncation. Thus $I$ is an ordinary
$U$-intertwining operator. If $S\subset R|_U$ is any simple submodule,
injectivity of $f_I$ implies that $I|_{S\times T}$ is nonzero.
Rigidity over $U$ then gives a nonzero, hence injective, morphism
\[
 S\longrightarrow T\boxtimes_U T^\vee.
\]
The conjugate-module description of the contragredient gives a unitary
form on $T^\vee$. Theorem~\ref{thm:positive-fusion}, applied to $U$,
makes $T\boxtimes_U T^\vee$ unitary. Restricting its positive form
unitarizes $S$. Taking the finite orthogonal sum over a simple
decomposition of $R|_U$ gives a positive $U$-invariant form $h_1$ on
$R$.

The form $h_1$ gives the correct adjoints for fields in $U$. It does
not yet give the correct adjoints for fields in $J$, although those
fields already act algebraically on $R$. We will construct a second
positive form from their action and then add the two forms.
Dong and Lin constructed unitary order-two simple-current VOA extensions
under compatible antiunitary-involution hypotheses
\cite[Theorem~3.3]{DongLin}.
Carpi, Gaudio, Giorgetti and Hillier proved unitarity and complete
unitarity of simple CFT-type extensions of a completely unitary base
\cite[Theorem~4.7]{CGGH}. Here the extension algebra $D$ is already
unitary, while the arbitrary module $R$ initially has a positive form
only after restriction to $U$. The argument below transfers that form
to the existing $D$-action.

\smallskip
\noindent\emph{A second form from the sign action.}
Since $V\ne\mathbb C$, the flip is faithful and $J\ne0$: if $v$ is
independent of $\Omega$, then $v\otimes\Omega-\Omega\otimes v$ is a
nonzero vector of $J$. The orbifold constituent theorem
\cite[Corollary~4.8]{McRaeOrbifold} identifies the tensor subcategory
generated by the constituents of $D$ with
$\operatorname{Rep}(\mathbb Z/2\mathbb Z)$. Here $D$ is an ordinary
VOA, so the abelian associativity cocycle in that theorem is trivial.
In particular, $J$ is a simple current, that is, a simple invertible
object, and
\[
 J\boxtimes_U J\cong U.
\]
The multiplication morphism
\[
 m:J\boxtimes_U J\longrightarrow U
\]
coming from $Y^D|_{J\times J}$ is nonzero. For example, the invariant
bilinear form of $D$ restricts nondegenerately to $J$, so some product
of two states in $J$ has a nonzero vacuum coefficient. Thus $m$ is an
isomorphism.

Let
\[
 L:J\boxtimes_U R\longrightarrow R
\]
be the morphism corresponding to the action $Y^R|_{J\times R}$.
Ordinary $D$-module associativity, expressed using the fusing
isomorphism over $U$, says that
\begin{equation}\label{eq:sign-action-square}
 L(1_J\boxtimes L)\alpha_{J,J,R}
   =\ell_R(m\boxtimes1_R).
\end{equation}
Here $\alpha_{J,J,R}:(J\boxtimes_U J)\boxtimes_U R
\longrightarrow J\boxtimes_U(J\boxtimes_U R)$ is the associator and
$\ell_R:U\boxtimes_U R\longrightarrow R$ is the left unit
isomorphism; undecorated fusion signs in this identity are over $U$.
The right side is an isomorphism, so $L$ is epic. Tensoring by the
invertible object $J$ preserves composition length. The two sides of
$L$ have the same finite length, and hence $L$ is an isomorphism.

Give $J$ the form inherited from $D$. Apply
Theorem~\ref{thm:positive-fusion} to $J$ and $(R,h_1)$, and transfer
its positive invariant fusion form along $L$. This gives a second
positive $U$-invariant form $h_2$ on $R$, for which
$Y^R|_{J\times R}$ is the transported universal intertwiner. We will
show that $h_1+h_2$ is invariant for all of $D$.

\smallskip
\noindent\emph{Adjunction exchanges the two forms.}
The creation intertwiner of the $U$-module $J$ is
$Y^D|_{J\times U}$ by skew symmetry. With the inherited forms, its
adjoint annihilation is
\begin{equation}\label{eq:sign-annihilation}
 \mathcal Y_J^-(\bar j,z)j_0
       =Y^D(\Theta^D j,z)j_0,
 \qquad j,j_0\in J.
\end{equation}
This follows from unitarity of $D$ and fixes the scalar normalization
of the annihilation. The linear identification $\bar J\cong J$ in
\eqref{eq:sign-annihilation} sends $\bar j$ to $\Theta^D j$.

For $w>0$, define the antilinear charge transform on $D$ by
\[
 t_w=e^{wL_1}(-w^{-2})^{L_0}\Theta^D.
\]
The integer grading of $D$ makes this expression unambiguous. For
$j,j'\in J$ and finite-energy $r_1,r_2\in R$, the exact transport
identity of Theorem~\ref{thm:positive-fusion}, followed by ordinary
$D$-module associativity, gives
\begin{equation}\label{eq:two-form-transport}
\begin{aligned}
 &h_2\bigl(Y^R(j,z)r_1,Y^R(j',w)r_2\bigr)\\
 &\quad=h_1\bigl(
 Y^R(Y^D(t_wj',w^{-1}-z)j,z)r_1,r_2\bigr)\\
 &\quad=h_1\bigl(
 Y^R(t_wj',w^{-1})Y^R(j,z)r_1,r_2\bigr).
\end{aligned}
\end{equation}
Initially, $0<z,w<1$ and $w^{-1}>z>w^{-1}-z>0$. Pairings of completed
vectors on the first line are the absolutely convergent sums over
equal energies supplied by the transport identity. On the last line,
$h_1(\,\cdot\,,r_2)$ is a restricted-dual functional: $h_1$ is
$U$-invariant and $U$ has the same conformal vector as $D$. Thus
convergent module associativity applies without already assuming
$D$-invariance of $h_1$.

The vectors $j_{(n)}r$, with $j\in J$, $r\in R$ and $n\in\mathbb Z$,
span $R$. Their span is a $U$-submodule by the intertwiner commutator
formula, and is the image of the surjective morphism $L$.
Comparing Laurent coefficients in $z$ in
\eqref{eq:two-form-transport} therefore yields
\begin{equation}\label{eq:two-form-adjoint-one}
 h_2\bigl(x,Y^R(j,w)y\bigr)
 =h_1\bigl(Y^R(t_wj,w^{-1})x,y\bigr),
 \qquad x,y\in R,\quad j\in J.
\end{equation}
For finite-energy $x,y$, the two sides are finite Laurent polynomials
in $w$. Indeed, the forms pair equal energies, the states of $D$ have
integer weights, and $e^{wL_1}j$ is a finite sum. Consequently,
\eqref{eq:two-form-adjoint-one} extends from its nonempty initial
interval to every $w>0$.

The charge transform satisfies
\begin{equation}\label{eq:charge-reflection-involution}
 t_{w^{-1}}t_w=1.
\end{equation}
For completeness, $\Theta^D$ commutes with the Virasoro operators and
with the real integer-weight signs. Conjugating $e^{wL_1}$ by
$(-w^2)^{L_0}$ gives $e^{-w^{-1}L_1}$, so
\[
 t_{w^{-1}}t_w
 =e^{w^{-1}L_1}(-w^2)^{L_0}
   e^{wL_1}(-w^{-2})^{L_0}
 =1.
\]
Apply \eqref{eq:two-form-adjoint-one} at radius $w^{-1}$ with charge
$t_wj$, exchange the two vector slots, and take complex conjugates.
Using \eqref{eq:charge-reflection-involution} gives the reverse relation
\begin{equation}\label{eq:two-form-adjoint-two}
 h_1\bigl(x,Y^R(j,w)y\bigr)
 =h_2\bigl(Y^R(t_wj,w^{-1})x,y\bigr).
\end{equation}
Adding \eqref{eq:two-form-adjoint-one} and
\eqref{eq:two-form-adjoint-two} proves the invariant adjunction for
$h=h_1+h_2$ on states of $J$. It already holds on $U$. Since
$D=U\oplus J$, the form $h$ is $D$-invariant; it is positive definite
as a sum of positive definite forms. Only ordinary associativity and
the normalized transport identity were used in this construction.

Finally, choose a nonzero homogeneous vector $q\in P^\vee$ and put
\[
 h_P(p,p')=h(p\otimes q,p'\otimes q),\qquad p,p'\in P.
\]
This is positive definite. The fields of $V\otimes\Omega$ act on
$P\otimes\mathbb Cq$ as the fields of $V$ on $P$, and the PCT charge
transform on $v\otimes\Omega$ is the transform on $v$ tensored with
$\Omega$. The $D$-invariant adjunction therefore restricts to the
$V$-invariant adjunction for $h_P$. This unitarizes the arbitrary
simple module $P$. Finite direct sums of these forms prove the final
assertion.
\end{proof}
